% !TEX root = ../main.tex
\section{Upper Bound with Increasing \texorpdfstring{$s_k$}{s\_k}}
\label{sec:upper-bound}

\todo{Main upper bound within the same family $\calC_\theta$: constant parameters with $s=\Theta(1/\epsilon)$.}
% with s = Theta(1/epsilon) (e.g., mu = 0, alpha = Theta(epsilon)), giving
% O(epsilon^{-3}). After the proof, compare its threshold on epsilon*s with
% c_1 of Theorem~\ref{thm:lower-bound}. The increasing-s_k result below then
% becomes the accuracy-independent variant.

Throughout this section, we additionally impose the following assumption.
It is not part of the problem class \(\mathcal I(\mathbf C,\Delta)\) of
Definition~\ref{def:problem-class}.

\begin{assumption}[Level-bounded objective]
\label{ass:level-boundedness}
For every \(c\in\RR\), the sublevel set
\(\{\bx\in\RR^n:F(\bx)\le c\}\) is empty or compact.
\end{assumption}

\subsection{Algorithm with increasing \texorpdfstring{$s_k$}{s\_k}}
\label{sec:method}

This section uses the updates of Definition~\ref{def:splalm} with
iteration-dependent parameters \(s_k\), \(\eta_k\), \(\tau_k\), \(\alpha_k\),
and \(\beta_k\), and with an arbitrary initialization; see
Algorithm~\ref{alg:ddlalm}. By~\eqref{eq:reduced-al}--\eqref{eq:potential},
the first two updates read
\begin{equation}\label{eq:explicit-updates}
    \begin{aligned}
        \bx^{k+1}
        &=
        \bx^k-\eta_k
        \Bigl[
            \nabla f(\bx^k)
            +J_{\calA}(\bx^k)^\top\bq(\bx^k,\blmbd^k)
            +\frac{1}{\gamma}(\bx^k-\bz^k)
        \Bigr], \\
        \blmbd^{k+1}
        &=
        \blmbd^k+\tau_k
        \Bigl[
            \calA(\bx^{k+1})
            -\by(\bx^{k+1},\blmbd^k)
            -\frac{1}{s_k}(\blmbd^k-\bmu^k)
        \Bigr],
    \end{aligned}
\end{equation}
and the outputs \((\bx^{k+1},\by^{k+1},\bq^k)\) are those of
Definition~\ref{def:splalm}.

\begin{algorithm}[t]
    \caption{Smoothed-Perturbed Linearized Augmented Lagrangian Method (SPLALM)}
    \label{alg:ddlalm}
    \small
    \begin{algorithmic}[1]
        \STATE \textbf{Input:} Initial
        \(\bx^0,\bz^0\in\RR^n\) and
        \(\blmbd^0,\bmu^0\in\RR^m\); 
        Parameters: \(\rho>0\) and \(\gamma>0\);
        sequences: \(\{\eta_k\}\), \(\{\tau_k\}\), \(\{\alpha_k\}\), and
        \(\{\beta_k\}\);
        and iteration limit \(K\).
        \FOR{\(k=0,1,\ldots,K-1\)}
            \STATE  \(\bx^{k+1}=\bx^k- \eta_k \nabla_{\bx} \mathcal P_{s_k}(\bx^{k},\blmbd^{k};\bz^{k},\bmu^{k})\);
            \STATE  \(\blmbd^{k+1}=\blmbd^k+\tau_k
            \nabla_{\blmbd}\mathcal P_{s_k}
            (\bx^{k+1},\blmbd^{k};\bz^{k},\bmu^{k})\).
            \STATE 
            \(\bz^{k+1}=\bz^k+ \alpha_k(\bx^{k+1}-\bz^k)\).
            \STATE  \(\bmu^{k+1}=\bmu^k+ \beta_k(\blmbd^{k+1}-\bmu^k)\).
        \ENDFOR
    \end{algorithmic}
\end{algorithm}

We use fixed primal and dual stepsizes together with diminishing center
stepsizes. Choose \(\rho>0\) and \(\gamma>0\) satisfying
Proposition~\ref{prop:curvature-smoothness}. With \(L_x\),
\(\mu_\gamma\), and \(L_{d,0}\) defined in
Section~\ref{sec:complexity-analysis}, set \(\eta_k\equiv\eta\) and
\(\tau_k\equiv\tau\). An admissible choice is
\begin{equation}\label{eq:ddlalm-parameter-schedule}
\begin{aligned}
    s_k &= c_s(k+1)^{1/3},
    \qquad
    \eta = \frac{1}{2L_x},
    \qquad
    \tau =
    \left(
        2L_{d,0}
        +\frac{72L_xG_{\cal A}^2}{\mu_\gamma^2}
    \right)^{-1},
    \\
    \alpha_k &=
    \frac{\min\{s_0,\gamma\mu_\gamma\tau\}}{36s_k},
    \qquad
    \beta_k=\frac{\theta_\mu}{k+k_0},
\end{aligned}
\end{equation}
where \(c_s>0\), \(\theta_\mu\in(0,1]\), and \(k_0\ge1\) are constants
independent of \(K\). Consequently,
\[
    s_k=\Theta((k+1)^{1/3}),
    \qquad
    \alpha_k=\Theta((k+1)^{-1/3}),
    \qquad
    \beta_k=\Theta((k+1)^{-1}).
\]

\begin{theorem}
\label{thm:main}
Suppose that Assumptions~\ref{ass:level-boundedness},
\ref{ass:smoothness}, and~\ref{ass:h} hold.
Let the sequence generated by \Cref{alg:ddlalm} use
\eqref{eq:ddlalm-parameter-schedule}, with $\rho$ and $\gamma$ as specified
above. For $K\ge2$, choose a tail-best index
\[
    \kappa\in\argmin_{k\in\{\lceil K/2\rceil,\ldots,K-1\}}
    \left\{
        \|\Delta_{\bx}^{k+1}\|^2
        +\|\Delta_{\blmbd}^{k+1}\|^2
        +\alpha_k\|\bx^{k+1}-\bz^k\|^2
    \right\}.
\]
Then it holds that
\[
\max\left\{
\left[
    \left\|
        \nabla f({\bx}^{\kappa+1})
        +
        J_{\cal A}({\bx}^{\kappa+1})^\top{\bq}^{\kappa}
    \right\|^2
\right],
\;
\left[
    \left\|
        {\cal A}({\bx}^{\kappa+1})-{\by}^{\kappa+1}
    \right\|^2
\right]
\right\}
=
\mathcal O(K^{-2/3}),
\]
with \({\bq}^{\kappa}\in\partial h({\by}^{\kappa+1})\).
Consequently, given \(\epsilon>0\), the iteration complexity of
\Cref{alg:ddlalm} to obtain an \(\epsilon\)-KKT point
is of order \(\mathcal O(\epsilon^{-3})\).
\end{theorem}



\subsection{Complexity Analysis}
\label{sec:complexity-analysis}
\label{sec:basic-descent}
Let us define the constants
\begin{equation}
    L_{\calD,\bx}
    :=
    L_f + L_{\calA}L_h + \rho G_{\calA}^{2},
    \label{eq:x-weak-convexity-modulus}
\end{equation}
and 
\[
    \mu_\gamma
    :=
    \frac{1}{\gamma}-L_{\calD,\bx},
    \qquad
    L_x
    :=
    L_{\calD,\bx}+\frac{1}{\gamma},
    \qquad
    L_\lambda(s)
    :=
    \frac{1}{\rho}+\frac{1}{s},
\] where $L_f$, $L_{\calA}$, $L_h$, $\rho$, and $G_{\calA}$ defined as before. 

The following proposition summarizes the curvature and smoothness properties
of the smoothed perturbed potential function $\mathcal P_s$.
The proof is deferred to Appendix~\ref{app:curvature-smoothness}.

\begin{proposition}
\label{prop:curvature-smoothness}
Suppose that Assumptions~\ref{ass:smoothness} and~\ref{ass:h} hold, and let the smoothing parameter  
$ 0 < \gamma < \frac{1}{2L_{\calD,\bx}}$. Then, for all $s>0$, $\mathcal P_s$ is
$\mu_\gamma $-strongly convex in $\bx$ and
$1/s$-strongly concave in $\blmbd$. Moreover, for all
$\bx,\bx'\in\RR^n$,
$\blmbd,\blmbd'\in\RR^m$,
$\bz\in\RR^n$, and
$\bmu\in\RR^m$, it holds that 
\begin{align}
    &\bigl\|
    \nabla_{\bx}\mathcal{P}_s(\bx,\blmbd,\bz,\bmu)
    -
    \nabla_{\bx}\mathcal{P}_s(\bx',\blmbd',\bz,\bmu)
    \bigr\|
    \le
    L_x\|\bx-\bx'\|
    +
    G_{\calA}\|\blmbd-\blmbd'\|,
    \label{eq:payoff-x-smoothness}
    \\
    &\bigl\|
    \nabla_{\blmbd}\mathcal{P}_s(\bx,\blmbd,\bz,\bmu)
    -
    \nabla_{\blmbd}\mathcal{P}_s(\bx',\blmbd',\bz,\bmu)
    \bigr\|
    \le
    G_{\calA}\|\bx-\bx'\|
    +
    L_\lambda(s)\|\blmbd-\blmbd'\|
    \label{eq:payoff-lambda-smoothness}.
\end{align}
In particular, $\mathcal P_s$ is $L_x$-smooth in $\bx$ and $L_\lambda(s)$-smooth in $\blmbd$.
\end{proposition}
 
\subsubsection{Value Functions and Lyapunov Function}
\label{subsec:value-functions}

Fix $s>0$ and $(\bz,\bmu)$. By Proposition~\ref{prop:curvature-smoothness},
$\mathcal P_s(\cdot,\cdot;\bz,\bmu)$ is strongly convex in $\bx$ and strongly
concave in $\blmbd$. Hence the following value functions are well defined,
and the minimizer and maximizer below are unique:
\begin{align}
    d_s(\blmbd;\bz,\bmu)
    &:= \min_{\bx\in\RR^n}\mathcal P_s(\bx,\blmbd;\bz,\bmu),
    &
    \bx_s(\blmbd;\bz,\bmu)
    &:= \argmin_{\bx\in\RR^n}\mathcal P_s(\bx,\blmbd;\bz,\bmu),
    \label{eq:dual-value}\\
    p_s(\bz,\bmu)
    &:= \max_{\blmbd\in\RR^m} d_s(\blmbd;\bz,\bmu),
    &
    \blmbd_s^\star(\bz,\bmu)
    &:= \argmax_{\blmbd\in\RR^m} d_s(\blmbd;\bz,\bmu).
    \label{eq:saddle-value}
\end{align}
Strong convexity--strong concavity also yields strong duality,
\begin{equation}
    p_s(\bz,\bmu)
    = \min_{\bx\in\RR^n}\max_{\blmbd\in\RR^m}
      \mathcal P_s(\bx,\blmbd;\bz,\bmu),
    \label{eq:minimax-equality}
\end{equation}
and the unique saddle point is
$(\bx_s^\star(\bz,\bmu),\blmbd_s^\star(\bz,\bmu))$, where
\[
    \bx_s^\star(\bz,\bmu)
    :=\bx_s\bigl(\blmbd_s^\star(\bz,\bmu);\bz,\bmu\bigr).
\]

\blue{Following the three-value Lyapunov construction for smoothed nonconvex--concave minimax methods~\cite{zhang2020smoothed}, and its perturbed-smoothed extension~\cite{li2026smoothing}, we adapt the construction to the reduced augmented-Lagrangian setting and define}
\begin{equation}
\begin{aligned}
    \Psi_s(\bx,\blmbd;\bz,\bmu)
    &:={}
    \mathcal P_s(\bx,\blmbd;\bz,\bmu)
    -2d_s(\blmbd;\bz,\bmu)
    +2p_s(\bz,\bmu) \\
    &= \underbrace{\mathcal P_s(\bx,\blmbd;\bz,\bmu)-d_s(\blmbd;\bz,\bmu)}_{\text{primal gap at fixed }\blmbd}
     + \underbrace{p_s(\bz,\bmu)-d_s(\blmbd;\bz,\bmu)}_{\text{dual gap}}
     + p_s(\bz,\bmu).
\end{aligned}
\label{eq:lyapunov-decomposition}
\end{equation}
Both gaps are nonnegative, so $\Psi_s\ge p_s$.

\paragraph{Iteration-$k$ shorthand.}
At iteration $k$ we evaluate the above objects at $(s,\bmu)=(s_k,\bmu^k)$,
replace the subscript $s$ by $k$, and suppress $\bmu^k$. For example,
$\mathcal P_k(\bx,\blmbd;\bz):=\mathcal P_{s_k}(\bx,\blmbd;\bz,\bmu^k)$ and
$d_k(\blmbd;\bz):=d_{s_k}(\blmbd;\bz,\bmu^k)$; the quantities
$p_k$, $\bx_k$, $\blmbd_k^\star$, $\bx_k^\star$, and $\Psi_k$ are defined
analogously. We also write
$\Delta_{\bx}^{k+1}:=\bx^{k+1}-\bx^k$, and define
$\Delta_{\blmbd}^{k+1}$, $\Delta_{\bz}^{k+1}$, and $\Delta_{\bmu}^{k+1}$
in the same way.

\subsubsection{Basic One-Step Estimate}
\label{subsec:basic-one-step}

Throughout this subsection, we assume that 
\begin{equation}
    s_{k+1}\ge s_k>0,\qquad
    0 < \tau_k\le s_k,\qquad
    0<\beta_k\le1,\qquad
    0\le\alpha_k\le2.
    \label{eq:basic-parameter-condition}
\end{equation} 
The following lemma gives a basic one-step descent estimate of the potential function $\mathcal P$ from iteration $k$ to $k+1$.

\begin{lemma}
\label{lem:basic-one-step}
Suppose Assumptions~\ref{ass:smoothness} and~\ref{ass:h} hold.
Let the iterates be generated by Algorithm~\ref{alg:ddlalm} with the
parameters satisfying \eqref{eq:basic-parameter-condition}.  Then, for
every $k\ge0$, the following holds:
\begin{equation}
\label{eq:basic-one-step}
\begin{aligned}
    &\mathcal P_{k+1}(\bx^{k+1},\blmbd^{k+1};\bz^{k+1})
    -\mathcal P_k(\bx^k,\blmbd^k;\bz^k) \\
    \le\;&
    -\left(\frac{1}{\eta_k}-\frac{L_x}{2}\right)
      \|\Delta_{\bx}^{k+1}\|^2
    +\left(\frac{1}{\tau_k}-\frac{1}{2s_k}\right)
      \|\Delta_{\blmbd}^{k+1}\|^2 \\
    &-\frac{\alpha_k(2-\alpha_k)}{2\gamma}
      \|\bx^{k+1}-\bz^k\|^2
    +\frac{\beta_k(2-\beta_k)}{2s_k}
      \|\blmbd^{k+1}-\bmu^k\|^2 +\frac{1}{2}
      \left(\frac{1}{s_k}-\frac{1}{s_{k+1}}\right)
      \|\blmbd^{k+1}-\bmu^{k+1}\|^2 .
\end{aligned}
\end{equation}
\end{lemma}


As shown in Lemma~\ref{lem:basic-one-step}, the positive ascent term
generated by the $\blmbd$-update prevents the potential
$\mathcal P_k$ from yielding a direct descent estimate. We therefore relate
this term to the one-step increase of the dual value function, as developed
in Lemma~\ref{lem:dual-value-gain}.  This relation allows us to construct a Lyapunov function with a descent property. Before proceeding, we establish several regularity properties of the value functions \(d_k\) and \(p_k\), with their gradient formulas following from Danskin's theorem.
\begin{lemma}
\label{lem:v2-value-regularity}
Under the conditions of Proposition~\ref{prop:curvature-smoothness}, for all $k\ge0$, the following properties hold:
\begin{enumerate}[label=(\roman*)]
\item
For every \(\bz \in \RR^n \), the minimizer
\(\bx_k(\blmbd;\bz)\) is unique and
\(\frac{G_{\cal A}}{\mu_\gamma}\)-Lipschitz continuous in
\(\blmbd\).

\item for every \(\bz\in\RR^n\),
\(d_k(\cdot;\bz)\) is \(1/s_k\)-strongly concave and
\(L_{d,k}\)-smooth, with gradient
\[
    \nabla_{\blmbd}d_k(\blmbd;\bz)
    =
    \nabla_{\blmbd}\mathcal P_k
    \bigl(\bx_k(\blmbd;\bz),\blmbd;\bz\bigr),
\]
where \(L_{d,k}:=\frac1\rho+\frac1{s_k}+\frac{G_{\cal A}^2}{\mu_\gamma}\).

\item
The function \(p_k(z)\) is continuously differentiable, with gradient
\[
    \nabla p_k(\bz)
    =
    \frac1\gamma
    \bigl(\bz-\bx_k^\star(\bz)\bigr),
\]
and has a \(1/\gamma\)-Lipschitz continuous gradient.
\end{enumerate}
\end{lemma}

The following lemma controls the dual ascent term.
\begin{lemma}
\label{lem:dual-value-gain}
Suppose that the conditions of
Lemma~\ref{lem:v2-value-regularity} hold, and let the iterates
be generated by Algorithm~\ref{alg:ddlalm} with parameters satisfying
\eqref{eq:basic-parameter-condition}.
For every $k\ge 0$, assume that
\begin{equation}
\label{eq:dual-gain-condition}
    \frac{1}{\tau_k}+\frac{1}{2s_k} \ge 2L_{d,k}.
\end{equation}
Then the following holds
\begin{equation*}
\begin{aligned}
& \left(\frac{1}{\tau_k}-\frac{1}{2s_k}\right)
  \|\Delta_{\boldsymbol{\lambda}}^{k+1}\|^2 \\
& \le 2\!\bigl[
    d_k(\boldsymbol{\lambda}^{k+1};\mathbf{z}^k)
    -d_k(\boldsymbol{\lambda}^k;\mathbf{z}^k)
  \bigr]
  -\frac{L_{d,k}}{2}
    \|\Delta_{\boldsymbol{\lambda}}^{k+1}\|^2
  +\frac{4G_{\mathcal A}^2}{L_{d,k}\mu_\gamma}
  \left[
    \mathcal P_k(
      \mathbf{x}^{k+1},\boldsymbol{\lambda}^k;\mathbf{z}^k)
    -d_k(\boldsymbol{\lambda}^k;\mathbf{z}^k)
  \right].
\end{aligned}
\end{equation*}
Moreover, it holds that
\begin{equation}
\begin{aligned}
&\mathcal P_{k+1}(\bx^{k+1},\blmbd^{k+1};\bz^{k+1})
-\mathcal P_k(\bx^k,\blmbd^k;\bz^k)
\\
&\le
-\left[
\frac1{\eta_k}-\frac{L_x}{2}
-
\frac{2G_{\mathcal A}^2(1/\eta_k+L_x)^2}
{L_{d,k}\mu_\gamma^2}
\right]
\|\Delta_{\bx}^{k+1}\|^2
-\frac{L_{d,k}}2\|\Delta_{\blmbd}^{k+1}\|^2 \\
&\qquad
+2\bigl[d_k(\blmbd^{k+1};\bz^k)-d_k(\blmbd^k;\bz^k)\bigr]
-\frac{\alpha_k(2-\alpha_k)}{2\gamma}
\|\bx^{k+1}-\bz^k\|^2 \\
&\qquad
+\frac{\beta_k(2-\beta_k)}{2s_k}
\|\blmbd^{k+1}-\bmu^k\|^2
+\frac12\left(\frac1{s_k}-\frac1{s_{k+1}}\right)
\|\blmbd^{k+1}-\bmu^{k+1}\|^2.
\end{aligned}
\label{eq:dual-value-full-step}
\end{equation}
\end{lemma}

The potential $\mathcal P_k$ increases along the dual update, so its variation alone does not capture progress in
both variables. We therefore seek a merit function that accounts for both
primal suboptimality and progress in dual maximization. To account for
these two effects, consider the partial potential
\begin{equation}
\label{eq:partial-potential}
    \mathcal P_k-cd_k
    =(\mathcal P_k-d_k)-(c-1)d_k,
    \qquad c>1.
\end{equation}
The first term measures the primal optimality gap, while the second term
decreases as the dual value increases. In the following analysis, we set
$c=2$, consistent with the coefficient of the dual-value increment in
\eqref{eq:dual-value-full-step}.

\begin{lemma}
\label{lem:partial-lyapunov-z-step}
Under the conditions of Lemma~\ref{lem:dual-value-gain}, for every $k\ge0$, the following holds
\begin{equation}
\begin{aligned}
&\Big[
\mathcal P_{k+1}(\bx^{k+1},\blmbd^{k+1};\bz^{k+1})
-2d_k(\blmbd^{k+1};\bz^{k+1})
\Big] -
\Big[
\mathcal P_k(\bx^k,\blmbd^k;\bz^k)
-2d_k(\blmbd^k;\bz^k)
\Big]\\
\le\,&
-\left[
\frac1{\eta_k}-\frac{L_x}{2}
-\frac{2G_{\mathcal A}^2(1/\eta_k+L_x)^2}
{L_{d,k}\mu_\gamma^2}
\right]
\|\Delta_{\bx}^{k+1}\|^2 -\frac{L_{d,k}}{2}\|\Delta_{\blmbd}^{k+1}\|^2 \\
&\qquad
+\frac{\alpha_k(2-\alpha_k)}{2\gamma}
\|\bx^{k+1}-\bz^k\|^2
-\frac{2\alpha_k}{\gamma}
\left\langle
\bx^{k+1}
-\bx_k(\blmbd^{k+1};\bz^{k+1}),
\bx^{k+1}-\bz^k
\right\rangle  \\
&\qquad
+\frac{\beta_k(2-\beta_k)}{2s_k}
\|\blmbd^{k+1}-\bmu^k\|^2
+\frac12\left(\frac1{s_k}-\frac1{s_{k+1}}\right)
\|\blmbd^{k+1}-\bmu^{k+1}\|^2.
\end{aligned}
\label{eq:partial-lyapunov-z-step}
\end{equation}
\end{lemma}
Equation~\eqref{eq:partial-lyapunov-z-step} gives the one-step estimate
for the partial potential in \eqref{eq:partial-potential} with $c=2$.
The $\bx$- and
$\blmbd$-updates contribute the first two terms on the right-hand side.
The $\bz$-update, however, produces a positive quadratic term and an
indefinite inner-product term. Hence, the partial potential does not yet
provide a sufficient descent estimate for the full iteration.

At an exact saddle point of
\(\min_{\bx}\max_{\blmbd}\mathcal P_k(\bx,\blmbd;\bz)\),
strong duality gives
\[
    \mathcal P_k\bigl(\bx_k^\star(\bz),
    \blmbd_k^\star(\bz);\bz\bigr)
    =d_k\bigl(\blmbd_k^\star(\bz);\bz\bigr)
    =p_k(\bz),
\]
and therefore
\[
    \mathcal P_k\bigl(\bx_k^\star(\bz),
    \blmbd_k^\star(\bz);\bz\bigr)
    -2d_k\bigl(\blmbd_k^\star(\bz);\bz\bigr)
    =-p_k(\bz).
\]
Consequently, any $\bz$-update that decreases the saddle value $p_k$
increases the partial potential $\mathcal P_k-2d_k$. This motivates us
to examine whether an additional value-function term can restore a descent
estimate.

We therefore add the value-function increment to the partial potential in
\eqref{eq:partial-potential}:
$2[p_k(\bz^{k+1})-p_k(\bz^k)]$ and combine the two estimates directly.
The resulting one-step bound is stated next.

\begin{lemma}
\label{lem:z-step-p}
Under the conditions of Lemma~\ref{lem:v2-value-regularity} and
$0\le\alpha_k\le2$, the $\bz$-update satisfies
\begin{equation}
\begin{aligned}
p_k(\bz^{k+1})-p_k(\bz^k)
\le\,
-\frac{\alpha_k(2-\alpha_k)}{2\gamma}
\|\bx^{k+1}-\bz^k\|^2 +
\frac{\alpha_k}{\gamma}
\left\langle
\bx^{k+1}-\bx_k^\star(\bz^k),
\bx^{k+1}-\bz^k
\right\rangle.
\end{aligned}
\label{eq:z-step-p}
% \notag
\end{equation}
\end{lemma}

\begin{lemma}
\label{lem:augmented-one-step}
Under the conditions of Lemma~\ref{lem:partial-lyapunov-z-step}, for every
$k\ge0$,
\begin{equation}
\begin{aligned}
&\Big[
\mathcal P_{k+1}(\bx^{k+1},\blmbd^{k+1};\bz^{k+1})
-2d_k(\blmbd^{k+1};\bz^{k+1})
+2p_k(\bz^{k+1})
\Big] \\
&\quad-
\Big[
\mathcal P_k(\bx^k,\blmbd^k;\bz^k)
-2d_k(\blmbd^k;\bz^k)
+2p_k(\bz^k)
\Big] \\
&\le
-\left[
\frac1{\eta_k}-\frac{L_x}{2}
-\frac{2G_{\mathcal A}^2(1/\eta_k+L_x)^2}
{L_{d,k}\mu_\gamma^2}
\right]
\|\Delta_{\bx}^{k+1}\|^2
-\frac{L_{d,k}}{2}\|\Delta_{\blmbd}^{k+1}\|^2 \\
&\qquad
-\frac{\alpha_k(2-\alpha_k)}{2\gamma}
\|\bx^{k+1}-\bz^k\|^2
+\frac{2\alpha_k}{\gamma}
\left\langle
\bx_k(\blmbd^{k+1};\bz^{k+1})
-\bx_k^\star(\bz^k),
\bx^{k+1}-\bz^k
\right\rangle \\
&\qquad
+\frac{\beta_k(2-\beta_k)}{2s_k}
\|\blmbd^{k+1}-\bmu^k\|^2
+\frac12\left(\frac1{s_k}-\frac1{s_{k+1}}\right)
\|\blmbd^{k+1}-\bmu^{k+1}\|^2.
\end{aligned}
\label{eq:augmented-one-step}
\end{equation}
\end{lemma}

Two gaps remain before $\Phi_k$ admits a genuine descent inequality: the augmented estimate still carries an indefinite cross term and positive multiplier/center-increment terms, and its left-hand side compares $d_k, p_k$ at both endpoints rather than the true $k \to k+1$ potential difference, which requires $d_{k+1}, p_{k+1}$. 

For that, we establish boundedness of the multipliers directly, rather than assuming compactness a priori: both multipliers are controlled by the AL structure and the inner residual, not by projection onto a compact set.
\begin{lemma}
\label{lem:current-multiplier-boundedness}
Suppose Assumptions~\ref{ass:smoothness} and~\ref{ass:h} hold.
Let the iterates be generated by Algorithm~\ref{alg:ddlalm} with
\eqref{eq:basic-parameter-condition} and
\begin{equation}
\label{eq:multiplier-boundedness-condition}
    \tau_k\left(\frac1\rho+\frac1{s_k}\right)\le1.
\end{equation}
Define the constant independent of $k$
\[
    M:=\max\{\|\blmbd^0\|,\|\bmu^0\|,L_h\}.
\]
Then, for every $k\ge0$, the following holds
\[
    \|\blmbd^k\|\le M,
    \qquad
    \|\bmu^k\|\le M.
\]
Moreover, for every $k\ge0$ and every $\bz\in\RR^n$, it holds that
\[
    \|\blmbd_k^\star(\bz)\|\le M.
\]
\end{lemma}

% The augmented one-step estimate also contains the indefinite term
% \[
% \left\langle
% \bx_k(\blmbd^{k+1};\bz^{k+1})-\bx_k^\star(\bz^k),
% \bx^{k+1}-\bz^k
% \right\rangle.
% \]
% Its first factor combines two effects: the dual iterate may differ from
% the dual maximizer, and the saddle-point primal component may move when
% $\bz$ changes. The next lemma quantifies both effects. The first estimate
% relates the dual suboptimality gap $p_k-d_k$ to the corresponding primal
% distance, while the second estimate controls the sensitivity of the
% primal minimizers to the center variable.

\begin{lemma}
\label{lem:saddle-center-stability}
Under the conditions of Proposition~\ref{prop:curvature-smoothness}, for every $k\ge0$, $\blmbd\in\RR^m$, and $\bz\in\RR^n$, it holds that
\begin{equation}
\label{eq:current-saddle-stability}
\|\bx_k(\blmbd;\bz)-\bx_k^\star(\bz)\|^2
\le
\frac{2}{\mu_\gamma}
\bigl[p_k(\bz)-d_k(\blmbd;\bz)\bigr].
\end{equation}
For any $\bz,\bz'\in\RR^n$,
\begin{equation}
\label{eq:current-center-stability}
\max\bigl\{
\|\bx_k(\blmbd;\bz')-\bx_k(\blmbd;\bz)\|,
\|\bx_k^\star(\bz')-\bx_k^\star(\bz)\|
\bigr\}
\le\frac{1}{\gamma\mu_\gamma}\|\bz'-\bz\|.
\end{equation}
\end{lemma}

We now establish an approximate descent inequality for
\Cref{alg:ddlalm} with respect to the Lyapunov function defined in \eqref{eq:lyapunov-decomposition}. We assume 
the parameters are chosen as
\begin{equation}
\label{eq:parameter-choice}
\begin{gathered}
    0<\gamma<\frac{1}{2L_{\calD,\bx}},\qquad
    0<s_k\le s_{k+1},\qquad
    0<\beta_k\le1,\qquad
    0<\eta_k\le\frac{1}{2L_x},\\
    0<\tau_k\le\left(\max\left\{2L_{d,k},\,
        \frac{36G_{\mathcal A}^2}{\eta_k\mu_\gamma^2}\right\}\right)^{-1},\qquad
    0<\alpha_k\le\frac1{36}\min\left\{1,\,
        \frac{\gamma\mu_\gamma\tau_k}{s_k}\right\},
\end{gathered}
\end{equation}
where $L_{d,k}:=\rho^{-1}+s_k^{-1}+G_{\mathcal A}^2/\mu_\gamma$.
Since $L_{d,k}\ge\rho^{-1}+s_k^{-1}$, \eqref{eq:parameter-choice} implies the conditions of Proposition~\ref{prop:curvature-smoothness},
Lemma~\ref{lem:current-multiplier-boundedness},
\eqref{eq:basic-parameter-condition},
\eqref{eq:dual-gain-condition}, and
\eqref{eq:multiplier-boundedness-condition}.

\begin{lemma}
\label{lem:moving-three-value-descent}
\blue{Suppose that Assumptions~\ref{ass:level-boundedness},
\ref{ass:smoothness}, and~\ref{ass:h} hold, and let}
$\{(\bx^k,\blmbd^k,\bz^k,\bmu^k)\}$ be generated by \Cref{alg:ddlalm}
with parameters satisfying~\eqref{eq:parameter-choice}. Then for every
$k\ge0$:
\begin{enumerate}[label=(\alph*)]
    \item for all $\bx\in\RR^n$, $\blmbd\in\RR^m$, and $\bz\in\RR^n$,
    \begin{equation}
    \label{eq:current-lyapunov-lower-bound}
        \Psi_k(\bx,\blmbd;\bz)\ge p_k(\bz)
        \ge\inf_{\bx\in\RR^n}F(\bx)-\frac{(L_h+M)^2}{2s_0};
    \end{equation}
    \item the iterates satisfy
    \begin{equation}
    \label{eq:full-lyapunov-descent}
    \begin{aligned}
        &\Psi_{k+1}(\bx^{k+1},\blmbd^{k+1};\bz^{k+1})
         -\Psi_k(\bx^k,\blmbd^k;\bz^k)\\
        &\qquad\le
         -\frac1{4\eta_k}\|\Delta_{\bx}^{k+1}\|^2
         -\frac1{8\tau_k}\|\Delta_{\blmbd}^{k+1}\|^2
         -\frac{\alpha_k}{2\gamma}\|\bx^{k+1}-\bz^k\|^2
         +\delta_k,
    \end{aligned}
    \end{equation}
    where
    $\delta_k:=10M^2\bigl(s_k^{-1}-s_{k+1}^{-1}\bigr)+20M^2\beta_k/s_k$.
\end{enumerate}
\end{lemma}


% For $s_k=c_s(k+1)^{1/3}$ and
% $\beta_k=\theta_\mu/(k+k_0)$ with $k_0\ge1$, the total error satisfies
% \begin{equation}
% \label{eq:explicit-transport-budget}
% \sum_{k=0}^\infty
% \left[
% 10M^2\left|\frac1{s_{k+1}}-\frac1{s_k}\right|
% +20M^2\frac{\beta_k}{s_k}
% \right]
% \le\frac{(10+80\theta_\mu)M^2}{c_s}.
% \end{equation}

% The augmented estimate adds $2p_k$ to the partial potential in
% \eqref{eq:partial-potential}. It uses the $1/\gamma$-smoothness of $p_k$
% once. The remaining
% inner product is the only new $\bz$-direction error that must be controlled
% by the subsequent stability estimates.

% \begin{lemma}[Saddle-point stability]
% \label{lem:saddle-point-stability}
% Under the conditions of
% Proposition~\ref{prop:curvature-smoothness}, for each $k\ge 0$ and
% each fixed $\bz$, let
% \[
%     \blmbd_k^\star(\bz)
%     :=\argmax_{\blmbd\in\RR^m}d_k(\blmbd;\bz).
% \]
% Then
% \[
%     \bx_k^\star(\bz)
%     =
%     \bx_k\bigl(\blmbd_k^\star(\bz),\bz\bigr),
% \]
% and
% \begin{equation}
% \label{eq:saddle-primal-stability}
% \left\|
%     \bx_k(\blmbd;\bz)-\bx_k^\star(\bz)
% \right\|
% \le
% \frac{G_{\mathcal A}}{\mu_\gamma}
% \left\|
%     \blmbd-\blmbd_k^\star(\bz)
% \right\|.
% \end{equation}
% Moreover,
% \begin{equation}
% \label{eq:dual-gap-to-multiplier-distance}
% \left\|
%     \blmbd-\blmbd_k^\star(\bz)
% \right\|^2
% \le
% 2s_k
% \bigl[
%     p_k(\bz)-d_k(\blmbd;\bz)
% \bigr].
% \end{equation}
% Consequently,
% \begin{equation}
% \label{eq:saddle-point-stability-final}
% \left\|
%     \bx_k(\blmbd;\bz)-\bx_k^\star(\bz)
% \right\|^2
% \le
% \frac{2G_{\mathcal A}^2s_k}{\mu_\gamma^2}
% \bigl[
%     p_k(\bz)-d_k(\blmbd;\bz)
% \bigr].
% \end{equation}
% \end{lemma}

The proof is given in Appendix~\ref{app:curvature-smoothness}.

\appendix
\subsection{Proofs of the Auxiliary Results}
\label{app:curvature-smoothness}

\begin{proof}[Proof of the curvature part of Proposition~\ref{prop:curvature-smoothness}]
For fixed $\blmbd$, the reduced payoff satisfies
\[
    \nabla_{\bx}\mathcal D_\rho(\bx,\blmbd)
    =
    \nabla f(\bx)+J_{\calA}(\bx)^\top\bq(\bx,\blmbd),
\]
where
\[
    \bq(\bx,\blmbd)
    =
    \operatorname{prox}_{\rho h^*}
    \bigl(\blmbd+\rho\calA(\bx)\bigr).
\]

The nonexpansiveness of the proximal mapping, the bounded Jacobian, and
the Lipschitz Jacobian imply
\[
    \|\bq(\bx,\blmbd)-\bq(\bx',\blmbd)\|
    \le \rho G_{\calA}\|\bx-\bx'\|.
\]
Assumption~\ref{ass:h} gives
\(\|\bq(\bx,\blmbd)\|\le L_h\).  Therefore
\[
    \|\nabla_{\bx}\mathcal D_\rho(\bx,\blmbd)
    -\nabla_{\bx}\mathcal D_\rho(\bx',\blmbd)\|
    \le L_{\calD,\bx}\|\bx-\bx'\|.
\]

Thus $\mathcal D_\rho(\cdot,\blmbd)$ is
$L_{\calD,\bx}$-weakly convex.  Adding
$\frac{1}{2\gamma}\|\bx-\bz\|^2$ gives
$\mu_\gamma$-strong convexity in $\bx$.

For fixed $(\bx,\bz,\bmu)$, $\mathcal D_\rho(\bx,\cdot)$ is concave in
$\blmbd$ because it is the pointwise infimum of affine functions of
$\blmbd$.  The term
$-\frac{1}{2s}\|\blmbd-\bmu\|^2$ is $1/s$-strongly concave.  Hence
$\mathcal P_s$ is $1/s$-strongly concave in $\blmbd$.
\end{proof}

\begin{proof}[Proof of the smoothness part of Proposition~\ref{prop:curvature-smoothness}]
The $\bx$-block bound follows from the Lipschitz estimate for
$\nabla_{\bx}\mathcal D_\rho$ in the proof of
Proposition~\ref{prop:curvature-smoothness} and the quadratic term
$\frac{1}{2\gamma}\|\bx-\bz\|^2$.

For the $\blmbd$-block, the identity
\[
    \nabla_{\blmbd}\mathcal P_s(\bx,\blmbd;\bz,\bmu)
    =
    \frac{1}{\rho}
    \bigl(\bq(\bx,\blmbd)-\blmbd\bigr)
    -\frac{1}{s}(\blmbd-\bmu)
\]
and the nonexpansiveness of
$\operatorname{prox}_{\rho h^*}$ give
$L_\lambda(s)=1/\rho+1/s$.
The same estimate, together with
\[
    \|\calA(\bx)-\calA(\bx')\|
    \le G_{\calA}\|\bx-\bx'\|,
\]
gives both cross Lipschitz bounds with constant $G_{\calA}$.
\end{proof}

\begin{proof}[Proof of Lemma~\ref{lem:dual-value-gain}]
Set
\[
    g^k
    :=
    \nabla_{\blmbd}\mathcal P_k
    (\bx^{k+1},\blmbd^k,\bz^k),
    \qquad
    \Delta_{\blmbd}^{k+1}=\tau_k g^k,
\]
where the second identity follows from the dual update.
By the $L_{d,k}$-smoothness of $d_k(\cdot;\bz^k)$,
\[
\begin{aligned}
&d_k(\blmbd^{k+1};\bz^k)-d_k(\blmbd^k;\bz^k) \\
&\quad\ge
\left\langle
\nabla_{\blmbd}d_k(\blmbd^k;\bz^k),
\Delta_{\blmbd}^{k+1}
\right\rangle
-\frac{L_{d,k}}{2}
\|\Delta_{\blmbd}^{k+1}\|^2.
\end{aligned}
\]
Danskin's theorem and Lemma~\ref{lem:v2-value-regularity} give
\[
\nabla_{\blmbd}d_k(\blmbd^k;\bz^k)
=
\nabla_{\blmbd}\mathcal P_k
\bigl(\bx_k(\blmbd^k;\bz^k),\blmbd^k,\bz^k\bigr).
\]
The cross-Lipschitz estimate in
Proposition~\ref{prop:curvature-smoothness} therefore yields
\[
\begin{aligned}
&\left\langle
\nabla_{\blmbd}d_k(\blmbd^k;\bz^k),
\Delta_{\blmbd}^{k+1}
\right\rangle \\
&\quad=
\tau_k\left\langle
\nabla_{\blmbd}\mathcal P_k
\bigl(\bx_k(\blmbd^k;\bz^k),\blmbd^k,\bz^k\bigr),g^k
\right\rangle \\
&\quad\ge
\tau_k\|g^k\|^2
-\tau_kG_{\cal A}
\|\bx^{k+1}-\bx_k(\blmbd^k;\bz^k)\|\|g^k\|.
\end{aligned}
\]
Using
\[
    \tau_k\|g^k\|^2
    =\frac{1}{\tau_k}\|\Delta_{\blmbd}^{k+1}\|^2,
    \qquad
    \tau_k\|g^k\|
    =\|\Delta_{\blmbd}^{k+1}\|,
\]
we obtain
\begin{equation}
\begin{aligned}
&d_k(\blmbd^{k+1};\bz^k)-d_k(\blmbd^k;\bz^k) \\
&\quad\ge
\left(\frac{1}{\tau_k}-\frac{L_{d,k}}{2}\right)
\|\Delta_{\blmbd}^{k+1}\|^2
-G_{\cal A}\|\bx^{k+1}-\bx_k(\blmbd^k;\bz^k)\|
\|\Delta_{\blmbd}^{k+1}\|.
\end{aligned}
\label{eq:dual-value-gain-appendix}
\end{equation}

Let
\[
    b_k:=
    \frac{1}{\tau_k}+\frac{1}{2s_k}-L_{d,k}.
\]
Multiplying \eqref{eq:dual-value-gain-appendix} by $2$ and rearranging gives
\[
\begin{aligned}
&\left(\frac{1}{\tau_k}-\frac{1}{2s_k}\right)
\|\Delta_{\blmbd}^{k+1}\|^2 \\
&\quad\le
2\bigl[d_k(\blmbd^{k+1};\bz^k)-d_k(\blmbd^k;\bz^k)\bigr]
-b_k\|\Delta_{\blmbd}^{k+1}\|^2 \\
&\qquad+
2G_{\cal A}\|\bx^{k+1}-\bx_k(\blmbd^k;\bz^k)\|
\|\Delta_{\blmbd}^{k+1}\|.
\end{aligned}
\]
The assumed parameter condition implies $b_k\ge L_{d,k}>0$. Young's inequality
with parameter $L_{d,k}$ gives
\[
2G_{\cal A}\|\bx^{k+1}-\bx_k(\blmbd^k;\bz^k)\|
\|\Delta_{\blmbd}^{k+1}\|
\le
\frac{L_{d,k}}{2}\|\Delta_{\blmbd}^{k+1}\|^2
+\frac{2G_{\cal A}^2}{L_{d,k}}
\|\bx^{k+1}-\bx_k(\blmbd^k;\bz^k)\|^2.
\]
Since $b_k\ge L_{d,k}$, we have
\[
    -b_k\|\Delta_{\blmbd}^{k+1}\|^2
    +\frac{L_{d,k}}{2}\|\Delta_{\blmbd}^{k+1}\|^2
    \le
    -\frac{L_{d,k}}{2}\|\Delta_{\blmbd}^{k+1}\|^2.
\]
Finally, the $\mu_\gamma$-strong convexity of
$\mathcal P_k(\cdot,\blmbd^k;\bz^k)$ gives
\[
\|\bx^{k+1}-\bx_k(\blmbd^k;\bz^k)\|^2
\le
\frac{2}{\mu_\gamma}
\left[
\mathcal P_k(\bx^{k+1},\blmbd^k;\bz^k)
-d_k(\blmbd^k;\bz^k)
\right].
\]
Substituting this estimate proves the first inequality.
\end{proof}

\begin{proof}[Proof of the second estimate in Lemma~\ref{lem:dual-value-gain}]
Fix $k\ge0$. By Proposition~\ref{prop:curvature-smoothness},
$\mathcal P_k(\cdot,\blmbd^k;\bz^k)$ is
$\mu_\gamma$-strongly convex. Applying the strong-convexity inequality
at $\bx^{k+1}$ to the minimizer $\bx_k(\blmbd^k;\bz^k)$ gives
\[
\begin{aligned}
&\mathcal P_k(\bx^{k+1},\blmbd^k;\bz^k)-d_k(\blmbd^k;\bz^k) \\
&\quad\le
\left\langle
\nabla_{\bx}\mathcal P_k(\bx^{k+1},\blmbd^k;\bz^k),
\bx^{k+1}-\bx_k(\blmbd^k;\bz^k)
\right\rangle \\
&\qquad
-\frac{\mu_\gamma}{2}
\|\bx^{k+1}-\bx_k(\blmbd^k;\bz^k)\|^2 \\
&\quad\le
\frac{1}{2\mu_\gamma}
\|\nabla_{\bx}\mathcal P_k(\bx^{k+1},\blmbd^k;\bz^k)\|^2,
\end{aligned}
\]
where the last inequality follows from Young's inequality.
The $\bx$-update and the $L_x$-Lipschitz continuity of
$\nabla_{\bx}\mathcal P_k(\cdot,\blmbd^k;\bz^k)$ imply
\[
\begin{aligned}
&\|\nabla_{\bx}\mathcal P_k(\bx^{k+1},\blmbd^k;\bz^k)\| \\
&\quad\le
\|\nabla_{\bx}\mathcal P_k(\bx^k,\blmbd^k;\bz^k)\|
+L_x\|\Delta_{\bx}^{k+1}\| \\
&\quad=
\left(\frac{1}{\eta_k}+L_x\right)
\|\Delta_{\bx}^{k+1}\|.
\end{aligned}
\]
Combining these estimates yields
\begin{equation}
\label{eq:primal-gap-control}
\mathcal P_k(\bx^{k+1},\blmbd^k;\bz^k)
-d_k(\blmbd^k;\bz^k)
\le
\frac{(1/\eta_k+L_x)^2}{2\mu_\gamma}
\|\Delta_{\bx}^{k+1}\|^2.
\end{equation}
Substituting \eqref{eq:primal-gap-control} into the first inequality of
Lemma~\ref{lem:dual-value-gain} gives a bound on the dual ascent term,
where the coefficient of
$\|\Delta_{\bx}^{k+1}\|^2$ becomes
\[
\frac{4G_{\mathcal A}^2}{L_{d,k}\mu_\gamma}
\frac{(1/\eta_k+L_x)^2}{2\mu_\gamma}
=
\frac{2G_{\mathcal A}^2(1/\eta_k+L_x)^2}
{L_{d,k}\mu_\gamma^2}.
\]
Substitute the resulting bound on
$(1/\tau_k-1/(2s_k))\|\Delta_{\blmbd}^{k+1}\|^2$
into \eqref{eq:basic-one-step} and combine the coefficients of
$\|\Delta_{\bx}^{k+1}\|^2$. The $\bz$-step term and both parameter-drift
terms remain unchanged. This proves \eqref{eq:dual-value-full-step}.
\end{proof}

\begin{proof}[Proof of Lemma~\ref{lem:partial-lyapunov-z-step}]
Subtract
$2[d_k(\blmbd^{k+1};\bz^k)-d_k(\blmbd^k;\bz^k)]$
from both sides of \eqref{eq:dual-value-full-step}. To replace
$d_k(\blmbd^{k+1};\bz^k)$ on the left by
$d_k(\blmbd^{k+1};\bz^{k+1})$, add
$-2[d_k(\blmbd^{k+1};\bz^{k+1})-d_k(\blmbd^{k+1};\bz^k)]$
to both sides.
By the minimum-value definition,
\[
d_k(\blmbd^{k+1};\bz^k)
\le
\mathcal P_k\bigl(\bx_k(\blmbd^{k+1};\bz^{k+1}),
\blmbd^{k+1};\bz^k\bigr).
\]
Subtracting this inequality from the corresponding equality at
$\bz^{k+1}$ and expanding the quadratic terms yields
\[
\begin{aligned}
&d_k(\blmbd^{k+1};\bz^{k+1})-d_k(\blmbd^{k+1};\bz^k) \\
&\quad\ge
\frac1\gamma\left\langle
\bz^k-\bx_k(\blmbd^{k+1};\bz^{k+1}),
\Delta_{\bz}^{k+1}\right\rangle
+\frac1{2\gamma}\|\Delta_{\bz}^{k+1}\|^2 \\
&\quad=
-\frac{\alpha_k(2-\alpha_k)}{2\gamma}
\|\bx^{k+1}-\bz^k\|^2 \\
&\qquad+
\frac{\alpha_k}{\gamma}\left\langle
\bx^{k+1}-\bx_k(\blmbd^{k+1};\bz^{k+1}),
\bx^{k+1}-\bz^k\right\rangle.
\end{aligned}
\]
Multiply this inequality by $-2$ and substitute it into the preceding
estimate. Its squared-step term combines with the $\bz$-step term in
\eqref{eq:dual-value-full-step} to give
$\alpha_k(2-\alpha_k)\|\bx^{k+1}-\bz^k\|^2/(2\gamma)$.
Both parameter-drift terms remain unchanged, proving
\eqref{eq:partial-lyapunov-z-step}.
\end{proof}

\begin{proof}[Proof of Lemma~\ref{lem:z-step-p}]
By Lemma~\ref{lem:v2-value-regularity}, $p_k$ has a
$1/\gamma$-Lipschitz continuous gradient and
\[
    \nabla p_k(\bz^k)
    =
    \frac1\gamma\bigl(\bz^k-\bx_k^\star(\bz^k)\bigr).
\]
The $\bz$-update gives
\[
    \bz^{k+1}-\bz^k
    =\alpha_k(\bx^{k+1}-\bz^k).
\]
Applying the smoothness inequality for $p_k$ and expanding the inner
product yields \eqref{eq:z-step-p}.
\end{proof}

\begin{proof}[Proof of Lemma~\ref{lem:augmented-one-step}]
By Lemma~\ref{lem:z-step-p},
\begin{equation}
\begin{aligned}
2\bigl[p_k(\bz^{k+1})-p_k(\bz^k)\bigr]
\le{}&
-\frac{\alpha_k(2-\alpha_k)}{\gamma}
\|\bx^{k+1}-\bz^k\|^2 \\
&+
\frac{2\alpha_k}{\gamma}
\left\langle
\bx^{k+1}-\bx_k^\star(\bz^k),
\bx^{k+1}-\bz^k
\right\rangle.
\end{aligned}
\end{equation}
Adding this inequality to
\eqref{eq:partial-lyapunov-z-step}, the quadratic terms combine as
\[
\frac{\alpha_k(2-\alpha_k)}{2\gamma}
-\frac{\alpha_k(2-\alpha_k)}{\gamma}
=
-\frac{\alpha_k(2-\alpha_k)}{2\gamma}.
\]
The two inner products combine into
\[
\frac{2\alpha_k}{\gamma}
\left\langle
\bx_k(\blmbd^{k+1};\bz^{k+1})
-\bx_k^\star(\bz^k),
\bx^{k+1}-\bz^k
\right\rangle.
\]
Substitution into \eqref{eq:partial-lyapunov-z-step} proves
\eqref{eq:augmented-one-step}.
\end{proof}

\begin{proof}[Proof of Lemma~\ref{lem:current-multiplier-boundedness}]
The proximal optimality condition gives
$\bq(\bx,\blmbd)\in\partial h(\by(\bx,\blmbd))$, and hence
$\|\bq(\bx,\blmbd)\|\le L_h$. Substituting
$\calA(\bx)-\by(\bx,\blmbd)
=(\bq(\bx,\blmbd)-\blmbd)/\rho$ into the dual update yields
\[
\blmbd^{k+1}
=\left[1-\tau_k\left(\frac1\rho+\frac1{s_k}\right)\right]\blmbd^k
+\frac{\tau_k}{\rho}\bq(\bx^{k+1},\blmbd^k)
+\frac{\tau_k}{s_k}\bmu^k.
\]
The coefficients are nonnegative and sum to one. Together with
$\bmu^{k+1}=(1-\beta_k)\bmu^k+\beta_k\blmbd^{k+1}$, this proves
$\|\blmbd^k\|,\|\bmu^k\|\le M$ by induction from the definition of $M$.
The center update then gives
$\|\bmu^{k+1}-\bmu^k\|\le2M\beta_k$.

At the saddle point, the dual first-order condition gives
\[
\blmbd_k^\star(\bz)
=\frac{s_k}{\rho+s_k}
  \bq\bigl(\bx_k^\star(\bz),\blmbd_k^\star(\bz)\bigr)
+\frac{\rho}{\rho+s_k}\bmu^k.
\]
Since both vectors in this convex combination have norm at most $M$,
$\|\blmbd_k^\star(\bz)\|\le M$.
\end{proof}

\begin{proof}[Proof of Lemma~\ref{lem:saddle-center-stability}]
Fix $k$ and $\bz$. Strong convexity gives the unique minimizer
$\bx_k(\blmbd;\bz)$ for each $\blmbd$. The finite, smooth,
$1/s_k$-strongly concave function $d_k(\cdot;\bz)$ tends to $-\infty$
as $\|\blmbd\|\to\infty$, so it has a unique maximizer.
The minimax equality and strong convexity in $\bx$ identify the unique
saddle pair and give
$\bx_k^\star(\bz)=\bx_k(\blmbd_k^\star(\bz);\bz)$.

The maximizing property at $\bx_k^\star(\bz)$ and strong convexity at
fixed $\blmbd$ imply
\[
\begin{aligned}
p_k(\bz)-d_k(\blmbd;\bz)
&\ge
\mathcal P_k(\bx_k^\star(\bz),\blmbd;\bz)
-\mathcal P_k(\bx_k(\blmbd;\bz),\blmbd;\bz) \\
&\ge\frac{\mu_\gamma}{2}
\|\bx_k(\blmbd;\bz)-\bx_k^\star(\bz)\|^2.
\end{aligned}
\]
This proves \eqref{eq:current-saddle-stability}.
For the center sensitivity of $\bx_k(\blmbd;\bz)$, add the
strong-convexity inequalities at the minimizers for $\bz$ and $\bz'$.
The only center-dependent term is $\|\bx-\bz\|^2/(2\gamma)$, so
\[
\mu_\gamma\|\bx_k(\blmbd;\bz')-\bx_k(\blmbd;\bz)\|^2
\le\frac1\gamma\left\langle
\bz'-\bz,\bx_k(\blmbd;\bz')-\bx_k(\blmbd;\bz)\right\rangle.
\]
For the saddle primal points, use the same argument for
$\max_{\blmbd}\mathcal P_k(\cdot,\blmbd;\bz)$, whose strong-convexity
modulus is also $\mu_\gamma$. Cauchy--Schwarz gives
\eqref{eq:current-center-stability} in both cases.
The latter argument also gives
\[
\left\langle\bz'-\bz,\bx_k^\star(\bz')-\bx_k^\star(\bz)\right\rangle
\ge\gamma\mu_\gamma
\|\bx_k^\star(\bz')-\bx_k^\star(\bz)\|^2.
\]
Together with the Danskin formula for $\nabla p_k$ and
$2\gamma\mu_\gamma>1$, this verifies the smoothness bound used above:
\[
\begin{aligned}
\gamma^2\|\nabla p_k(\bz')-\nabla p_k(\bz)\|^2
&\le\|\bz'-\bz\|^2
-(2\gamma\mu_\gamma-1)
\|\bx_k^\star(\bz')-\bx_k^\star(\bz)\|^2 \\
&\le\|\bz'-\bz\|^2.
\end{aligned}
\]
\end{proof}

\begin{proof}[Proof of Lemma~\ref{lem:moving-three-value-descent}]
Let $u_{k+1}:=(\bx^{k+1},\blmbd^{k+1},\bz^{k+1})$ and define
\[
\widehat\Psi_{k+1|k}:=
\mathcal P_{k+1}(\bx^{k+1},\blmbd^{k+1};\bz^{k+1})
-2d_k(\blmbd^{k+1};\bz^{k+1})+2p_k(\bz^{k+1}).
\]
Then
\[
\Psi_{k+1}(u_{k+1})-\Psi_k(u_k)
=\widehat\Psi_{k+1|k}-\Psi_k(u_k)+R_k,
\]
where
\[
R_k:=2\bigl[d_k(\blmbd^{k+1};\bz^{k+1})
-d_{k+1}(\blmbd^{k+1};\bz^{k+1})\bigr]
+2\bigl[p_{k+1}(\bz^{k+1})-p_k(\bz^{k+1})\bigr].
\]

Apply Lemma~\ref{lem:augmented-one-step} to the first difference.
Decompose the indefinite term at $\bx_k^\star(\bz^{k+1})$ and apply
Lemma~\ref{lem:saddle-center-stability}, followed by Young's inequality.
The four bounds in \eqref{eq:parameter-choice} absorb the resulting
cross terms and leave
\[
\begin{aligned}
\widehat\Psi_{k+1|k}-\Psi_k(u_k)
\le{}&
-\frac1{4\eta_k}\|\Delta_{\bx}^{k+1}\|^2
-\frac1{8\tau_k}\|\Delta_{\blmbd}^{k+1}\|^2
-\frac{\alpha_k}{2\gamma}\|\bx^{k+1}-\bz^k\|^2 \\
&+\frac{\beta_k(2-\beta_k)}{2s_k}
\|\blmbd^{k+1}-\bmu^k\|^2
\\
&+\frac12\left(\frac1{s_k}-\frac1{s_{k+1}}\right)
\|\blmbd^{k+1}-\bmu^{k+1}\|^2.
\end{aligned}
\]

It remains to bound $R_k$. Lemma~\ref{lem:current-multiplier-boundedness}
gives $\|\blmbd^j\|,\|\bmu^j\|,\|\blmbd_j^\star(\bz)\|\le M$ and
$\|\bmu^{k+1}-\bmu^k\|\le2M\beta_k$. Direct comparison of the two
quadratic damping terms gives, uniformly in $\bx$ and $\bz$,
\[
|d_{k+1}(\blmbd;\bz)-d_k(\blmbd;\bz)|
\le\Delta_k,
\qquad
|p_{k+1}(\bz)-p_k(\bz)|\le\Delta_k,
\]
where
\[
\Delta_k:=2M^2\left(\frac1{s_k}-\frac1{s_{k+1}}\right)
+\frac{2M}{s_k}\|\bmu^{k+1}-\bmu^k\|.
\]
Consequently,
\[
|R_k|\le4\Delta_k
\le8M^2\left(\frac1{s_k}-\frac1{s_{k+1}}\right)
+16M^2\frac{\beta_k}{s_k}.
\]
The two explicit positive terms above satisfy
\[
\frac{\beta_k(2-\beta_k)}{2s_k}
\|\blmbd^{k+1}-\bmu^k\|^2
\le4M^2\frac{\beta_k}{s_k},
\]
and
\[
\frac12\left(\frac1{s_k}-\frac1{s_{k+1}}\right)
\|\blmbd^{k+1}-\bmu^{k+1}\|^2
\le2M^2\left(\frac1{s_k}-\frac1{s_{k+1}}\right).
\]
Adding the estimates proves \eqref{eq:full-lyapunov-descent}.

For the stated schedule, the two error series satisfy
\[
\sum_{k=0}^\infty\left(\frac1{s_k}-\frac1{s_{k+1}}\right)=\frac1{c_s},
\qquad
\sum_{k=0}^\infty\frac{\beta_k}{s_k}
\le\frac{4\theta_\mu}{c_s},
\]
which makes the accumulated transport error finite. The lower bound
\eqref{eq:current-lyapunov-lower-bound} follows as follows. For any
$\blmbd\in\partial h(\calA(\bx))$, proximal optimality gives
$\by(\bx,\blmbd)=\calA(\bx)$ and hence
$\mathcal D_\rho(\bx,\blmbd)=F(\bx)$. Since
$\|\blmbd\|\le L_h$ and $\|\bmu^k\|\le M$,
\[
\max_{\blmbd}\mathcal P_k(\bx,\blmbd;\bz)
\ge F(\bx)+\frac1{2\gamma}\|\bx-\bz\|^2
-\frac{(L_h+M)^2}{2s_k}.
\]
Minimizing in $\bx$ and using $s_k\ge s_0$ proves
\eqref{eq:current-lyapunov-lower-bound}, because
$\Psi_k=(\mathcal P_k-d_k)+(p_k-d_k)+p_k\ge p_k$.
Summing \eqref{eq:full-lyapunov-descent} and using the error budget
then gives
\[
\begin{aligned}
&\sum_{k=0}^{N-1}\left[
\frac1{4\eta_k}\|\Delta_{\bx}^{k+1}\|^2
+\frac1{8\tau_k}\|\Delta_{\blmbd}^{k+1}\|^2
+\frac{\alpha_k}{2\gamma}\|\bx^{k+1}-\bz^k\|^2\right] \\
&\quad\le
\Psi_0(\bx^0,\blmbd^0;\bz^0)-\inf_{\bx\in\RR^n}F(\bx)
+\frac{(L_h+M)^2}{2s_0}
+\frac{(10+80\theta_\mu)M^2}{c_s},
\qquad N\ge1.
\end{aligned}
\]
\end{proof}



\begin{proof}[Proof of Theorem~\ref{thm:main}]
The schedule \eqref{eq:ddlalm-parameter-schedule} satisfies
\eqref{eq:parameter-choice}: $\eta_kL_x=1/2$,
$\tau_kL_{d,k}\le\tau_kL_{d,0}\le1/2$,
$\tau_kG_{\cal A}^2\le\mu_\gamma^2/(72L_x)
=\eta_k\mu_\gamma^2/36$, and the two bounds on $\alpha_k$ follow
from $s_k\ge s_0$ and its definition. Moreover,
$s_{k+1}\ge s_k>0$, $\eta_k,\tau_k>0$, and
$0<\beta_k\le1$. Since $L_{d,k}\ge1/\rho+1/s_k$,
the multiplier-boundedness condition also holds.
By assumption, $\inf_{\bx}F(\bx)>-\infty$. The master budget proved after
Lemma~\ref{lem:moving-three-value-descent} is therefore finite and
independent of $K$.

The tail has $\lfloor K/2\rfloor\ge K/3$ elements. Since the primal
and dual stepsizes are fixed and
$\alpha_k\ge\alpha_{K-1}=\Theta(K^{-1/3})$ on the tail, the master
budget and the definition of $\kappa$ give
\[
\|\Delta_{\bx}^{\kappa+1}\|^2=O(K^{-1}),\qquad
\|\Delta_{\blmbd}^{\kappa+1}\|^2=O(K^{-1}),\qquad
\|\bx^{\kappa+1}-\bz^\kappa\|^2=O(K^{-2/3}).
\]
The output certificate follows directly from
\eqref{eq:prox-multiplier}. For stationarity, the $\bx$-update
and $L_{\calD,\bx}$-smoothness of $\mathcal D_\rho(\cdot,\blmbd^k)$ give
\[
\|\nabla f(\bx^{k+1})+J_{\cal A}(\bx^{k+1})^\top\bq^k\|
\le\left(\frac1{\eta_k}+L_x\right)\|\Delta_{\bx}^{k+1}\|
+\frac1\gamma\|\bx^{k+1}-\bz^k\|.
\]
For feasibility, the dual update gives
\[
{\cal A}(\bx^{k+1})-\by^{k+1}
=\frac1{\tau_k}\Delta_{\blmbd}^{k+1}
+\frac1{s_k}(\blmbd^k-\bmu^k),
\]
and hence
\[
\|{\cal A}(\bx^{k+1})-\by^{k+1}\|
\le\frac1{\tau_k}\|\Delta_{\blmbd}^{k+1}\|+\frac{2M}{s_k}.
\]
Square the last two bounds using $(a+b)^2\le2a^2+2b^2$ and use
$s_\kappa^{-2}=O(K^{-2/3})$. Both squared residuals are
$O(K^{-2/3})$, so their norms are $O(K^{-1/3})$. All constants depend
only on the problem data and fixed initialization and parameters, so
$K=O(\epsilon^{-3})$ suffices.
\end{proof}
